% GENERATED FILE. Edit canonical lessons, then run npm run generate:books.
% canonical-source-hash: fnv1a64:75add84b0c7a99e5
% canonical-lessons: GU-C421-ghoghat, GU-C421-hobalo, GU-C421-dhor, GU-C421-ghurakvu, GU-C421-bhagadvu

\chapter{Noise and Stray Cattle}
\label{ch:gu-noise-and-stray-cattle}
\hlchaptermodality{421}

\begin{chapteropening}
\textbf{By the end of this chapter:} \emph{I can complain about noise and animals in the street.}

Complete the last lesson of chapter 421: I can complain about noise and animals in the street.
\end{chapteropening}

\section[\texorpdfstring{\gu{ઘોંઘાટ} (ghõghāṭ)}{ઘોંઘાટ (ghõghāṭ)}]{\gu{ઘોંઘાટ} (ghõghāṭ) — noise; a din}

\label{lesson:GU-C421-ghoghat}



\begin{quote}
Before the new one: say the Gujarati for a post, then the Gujarati for to thatch.
\end{quote}

\subsection*{You'll want to know: \gu{ઘોંઘાટ}}
\textbf{\gu{ઘોંઘાટ}} — \emph{ghõghāṭ} — \textquotedblleft{}noise; a din\textquotedblright{}.

\begin{quote}
\gu{શહેરમાં} \gu{આખો} \gu{દિવસ} \gu{ઘોંઘાટ} \gu{રહે} \gu{છે}.
\end{quote}

\begin{quote}
\emph{śahermā̃ ākho divas ghõghāṭ rahe che.}
\end{quote}

\begin{quote}
In the city there is noise all day.
\end{quote}

A masculine noun.

\begin{cousinweb}
\textbf{\gu{ઘોંઘાટ}} (\emph{ghõghāṭ}) sounds like the noise it names. No dictionary we use records an older form, so its history is not settled.
\end{cousinweb}

\begin{grammarlens}[title={What you've built}]
One more word for getting around.
\end{grammarlens}

\subsection*{Your turn}
\begin{itemize}
\raggedright
  \item \emph{Say it:} \emph{ghõghāṭ}
  \item \emph{Say it:} \emph{ghõghāṭ}, once more
  \item \emph{Say it:} \emph{śahermā̃ ākho divas ghõghāṭ rahe che.} — In the city there is noise all day.
  \item \emph{Recall:} say the Gujarati for a post, then the Gujarati for to thatch, then say \emph{ghõghāṭ} again
\end{itemize}

\begin{tcolorbox}[breakable,colback=teal!4,colframe=teal!35!black,title={Before you move on}]
What does \gu{ઘોંઘાટ} mean? (\textquotedblleft{}Noise; a din\textquotedblright{}.) Say it once more.
\end{tcolorbox}
\section[\texorpdfstring{\gu{હોબાળો} (hobāḷo)}{હોબાળો (hobāḷo)}]{\gu{હોબાળો} (hobāḷo) — an uproar; an outcry}

\label{lesson:GU-C421-hobalo}



\begin{quote}
Before the new one: say the Gujarati for to thatch, then the Gujarati for a din.
\end{quote}

\subsection*{You'll want to know: \gu{હોબાળો}}
\textbf{\gu{હોબાળો}} — \emph{hobāḷo} — \textquotedblleft{}an uproar; an outcry\textquotedblright{}.

\begin{quote}
\gu{ભાવ} \gu{વધ્યા} \gu{એટલે} \gu{બજારમાં} \gu{હોબાળો} \gu{થયો}.
\end{quote}

\begin{quote}
\emph{bhāv vadhyā eṭle bajārmā̃ hobāḷo thayo.}
\end{quote}

\begin{quote}
Prices went up, so there was an uproar in the market.
\end{quote}

A masculine noun.

\begin{cousinweb}
\textbf{\gu{હોબાળો}} (\emph{hobāḷo}) No dictionary we use records an older form, so its history is not settled.
\end{cousinweb}

\begin{grammarlens}[title={What you've built}]
One more word for getting around.
\end{grammarlens}

\subsection*{Your turn}
\begin{itemize}
\raggedright
  \item \emph{Say it:} \emph{hobāḷo}
  \item \emph{Say it:} \emph{hobāḷo}, once more
  \item \emph{Say it:} \emph{bhāv vadhyā eṭle bajārmā̃ hobāḷo thayo.} — Prices went up, so there was an uproar in the market.
  \item \emph{Recall:} say the Gujarati for to thatch, then the Gujarati for a din, then say \emph{hobāḷo} again
\end{itemize}

\begin{tcolorbox}[breakable,colback=teal!4,colframe=teal!35!black,title={Before you move on}]
What does \gu{હોબાળો} mean? (\textquotedblleft{}An uproar; an outcry\textquotedblright{}.) Say it once more.
\end{tcolorbox}
\section[\texorpdfstring{\gu{ઢોર} (ḍhor)}{ઢોર (ḍhor)}]{\gu{ઢોર} (ḍhor) — cattle; livestock}

\label{lesson:GU-C421-dhor}



\begin{quote}
Before the new one: say the Gujarati for a din, then the Gujarati for an uproar.
\end{quote}

\subsection*{You'll want to know: \gu{ઢોર}}
\textbf{\gu{ઢોર}} — \emph{ḍhor} — \textquotedblleft{}cattle; livestock\textquotedblright{}.

\begin{quote}
\gu{રસ્તા} \gu{પર} \gu{ઢોર} \gu{બેઠાં} \gu{છે}, \gu{એટલે} \gu{ગાડીઓ} \gu{ધીમે} \gu{ચાલે} \gu{છે}.
\end{quote}

\begin{quote}
\emph{rastā par ḍhor beṭhā̃ che, eṭle gāḍīo dhīme cāle che.}
\end{quote}

\begin{quote}
Cattle are sitting on the road, so the cars move slowly.
\end{quote}

A neuter noun.

\begin{cousinweb}
\textbf{\gu{ઢોર}} (\emph{ḍhor}) CDIAL lists Hindi \emph{ḍhor}, \textquotedblleft{}cattle\textquotedblright{}, under Sanskrit \emph{dhaurēya}, \textquotedblleft{}a beast of burden\textquotedblright{} (CDIAL 6884), the word behind \gu{ધોરી} in \gu{ધોરીમાર્ગ}; it gives no Gujarati form for \gu{ઢોર}, so the link is hedged.
\end{cousinweb}

\begin{grammarlens}[title={What you've built}]
One more word for getting around.
\end{grammarlens}

\subsection*{Your turn}
\begin{itemize}
\raggedright
  \item \emph{Say it:} \emph{ḍhor}
  \item \emph{Say it:} \emph{ḍhor}, once more
  \item \emph{Say it:} \emph{rastā par ḍhor beṭhā̃ che, eṭle gāḍīo dhīme cāle che.} — Cattle are sitting on the road, so the cars move slowly.
  \item \emph{Recall:} say the Gujarati for a din, then the Gujarati for an uproar, then say \emph{ḍhor} again
\end{itemize}

\begin{tcolorbox}[breakable,colback=teal!4,colframe=teal!35!black,title={Before you move on}]
What does \gu{ઢોર} mean? (\textquotedblleft{}Cattle; livestock\textquotedblright{}.) Say it once more.
\end{tcolorbox}
\section[\texorpdfstring{\gu{ઘુરકવું} (ghurakvũ)}{ઘુરકવું (ghurakvũ)}]{\gu{ઘુરકવું} (ghurakvũ) — to growl; to snarl}

\label{lesson:GU-C421-ghurakvu}



\begin{quote}
Before the new one: say the Gujarati for an uproar, then the Gujarati for cattle.
\end{quote}

\subsection*{You'll want to know: \gu{ઘુરકવું}}
\textbf{\gu{ઘુરકવું}} — \emph{ghurakvũ} — \textquotedblleft{}to growl; to snarl\textquotedblright{}.

\begin{quote}
\gu{રાત્રે} \gu{શેરીમાં} \gu{કૂતરો} \gu{ઘુરકે} \gu{છે}.
\end{quote}

\begin{quote}
\emph{rātre śerīmā̃ kūtro ghurke che.}
\end{quote}

\begin{quote}
At night a dog growls in the street.
\end{quote}

Past (it …): \gu{તે} \gu{કૂતરો} \gu{ઘુરક્યો}, \emph{te kūtro ghurakyo}.

\begin{cousinweb}
\textbf{\gu{ઘુરકવું}} (\emph{ghurakvũ}) is a tadbhava of Sanskrit \emph{ghurati}, \textquotedblleft{}cries out frightfully\textquotedblright{}; CDIAL lists the Gujarati verb, \textquotedblleft{}to growl\textquotedblright{} (CDIAL 4487).
\end{cousinweb}

\begin{grammarlens}[title={What you've built}]
One more word for getting around.
\end{grammarlens}

\subsection*{Your turn}
\begin{itemize}
\raggedright
  \item \emph{Say it:} \emph{ghurakvũ}
  \item \emph{Say it:} \emph{ghurakvũ}, once more
  \item \emph{Say it:} \emph{rātre śerīmā̃ kūtro ghurke che.} — At night a dog growls in the street.
  \item \emph{Recall:} say the Gujarati for an uproar, then the Gujarati for cattle, then say \emph{ghurakvũ} again
\end{itemize}

\begin{tcolorbox}[breakable,colback=teal!4,colframe=teal!35!black,title={Before you move on}]
What does \gu{ઘુરકવું} mean? (\textquotedblleft{}To growl; to snarl\textquotedblright{}.) Say it once more.
\end{tcolorbox}
\section[\texorpdfstring{\gu{ભગાડવું} (bhagāḍvũ)}{ભગાડવું (bhagāḍvũ)}]{\gu{ભગાડવું} (bhagāḍvũ) — to chase away; to drive off}

\label{lesson:GU-C421-bhagadvu}



\begin{quote}
Before the new one: say the Gujarati for cattle, then the Gujarati for to growl.
\end{quote}

\subsection*{You'll want to know: \gu{ભગાડવું}}
\textbf{\gu{ભગાડવું}} — \emph{bhagāḍvũ} — \textquotedblleft{}to chase away; to drive off\textquotedblright{}.

\begin{quote}
\gu{ચોકીદાર} \gu{ઢોરને} \gu{દરવાજા} \gu{પાસેથી} \gu{ભગાડે} \gu{છે}.
\end{quote}

\begin{quote}
\emph{cokīdār ḍhorne darvājā pāsethī bhagāḍe che.}
\end{quote}

\begin{quote}
The watchman chases the cattle away from the gate.
\end{quote}

Past, with \gu{તેણે} (he or she …): \gu{તેણે} \gu{કૂતરાને} \gu{ભગાડ્યો}, \emph{teṇe kūtrāne bhagāḍyo}.

\begin{cousinweb}
\textbf{\gu{ભગાડવું}} (\emph{bhagāḍvũ}) is the causative partner of \gu{ભાગવું}, \textquotedblleft{}to run away\textquotedblright{}, which CDIAL lists under Sanskrit \emph{bhajyátē}, \textquotedblleft{}is broken\textquotedblright{} (CDIAL 9361): to make someone run off.
\end{cousinweb}

\begin{grammarlens}[title={What you've built}]
One more word for getting around.
\end{grammarlens}

\subsection*{Your turn}
\begin{itemize}
\raggedright
  \item \emph{Say it:} \emph{bhagāḍvũ}
  \item \emph{Say it:} \emph{bhagāḍvũ}, once more
  \item \emph{Say it:} \emph{cokīdār ḍhorne darvājā pāsethī bhagāḍe che.} — The watchman chases the cattle away from the gate.
  \item \emph{Recall:} say the Gujarati for cattle, then the Gujarati for to growl, then say \emph{bhagāḍvũ} again
\end{itemize}

\begin{tcolorbox}[breakable,colback=teal!4,colframe=teal!35!black,title={Before you move on}]
What does \gu{ભગાડવું} mean? (\textquotedblleft{}To chase away; to drive off\textquotedblright{}.) Say it once more.
\end{tcolorbox}
